% Emacs, this is -*-latex-*

\ex{Ortogonaliza\'c\~ao de Gram--Schmidt}
\label{ex:gram-schmidt} 

\begin{exenumerate}
    \item Dados dois vetores $\a_1$ e $\a_2$ em $\mathbb{R}^n$, mostre que 
    \begin{align}
        \u_1 &= \a_1 \\
        \u_2 &= \a_2 - \frac{\u_1^\top \a_2}{\u_1^\top \u_1}\u_1
    \end{align}
    s\~ao ortogonais entre si.
    \begin{solution}
        Dois vetores $\u_1$ e $\u_2$ de $\mathbb{R}^n$ s\~ao ortogonais
        se o seu produto interno for igual a zero. Calculando o produto interno $\u_1^\top \u_2$ obtemos
        \begin{align}
            \u_1^{\top}\u_2 &= \u_1^\top(\a_2 - \frac{\u_1^\top\a_2}{\u_1^\top\u_1}\u_1) \\
            &= \u_1^\top\a_2 - \frac{\u_1^\top\a_2}{\u_1^\top\u_1}\u_1^\top\u_1 \\
            &= \u_1^\top\a_2 - \u_1^\top\a_2 \\
            &= 0.
        \end{align}
        Portanto, os vetores $\u_1$ e $\u_2$ s\~ao ortogonais.
        
        Se $\a_2$ for um m\'ultiplo de $\a_1$, o procedimento de ortogonaliza\'c\~ao produz
        um vetor nulo para $\u_2$. Para ver isso, considere $\a_2 = \alpha \a_1$ para algum
        n\'umero real $\alpha$. Obtemos ent\~ao
        \begin{align} \u_2 &= \a_2 - \frac{\u_1^\top \a_2}{\u_1^\top \u_1}\u_1\\
            &= \alpha \u_1 - \frac{\alpha \u_1^\top \u_1}{\u_1^\top \u_1}\u_1\\
            &= \alpha \u_1 - \alpha \u_1\\ &=\mathbf{0}.
        \end{align}
    \end{solution}
    
    \item Mostre que qualquer combina\'c\~ao linear de $\a_1$ e $\a_2$ (linearmente independentes) pode ser escrita em termos de $\u_1$ e $\u_2$. 
    \begin{solution}
        Seja $\v$ uma combina\'c\~ao linear de $\a_1$ e $\a_2$, i.e. $\v = \alpha{\a_1} + \beta{\a_2}$ para alguns n\'umeros reais $\alpha$ e $\beta$.
        Expressando $\u_1$ e $\u_2$ em termos de $\a_1$ e $\a_2$, podemos escrever $\v$ como
        \begin{align}
            \v  &= \alpha \a_1 + \beta\a_2 \\
            &= \alpha\u_1 + \beta(\u_2 + \frac{\u_1^\top\a_2}{\u_1^\top\u_1}\u_1) \\
            &= \alpha\u_1 + \beta\u_2 + \beta \frac{\u_1^\top\a_2}{\u_1^\top\u_1}\u_1 \\
            &= (\alpha + \beta\frac{\u_1^\top\a_2}{\u_1^\top\u_1})\u_1 + \beta\u_2,
        \end{align}
        Como $\alpha + \beta((\u_1^\top\a_2)/(\u_1^\top\u_1))$ e $\beta$ s\~ao n\'umeros
        reais, podemos escrever $\v$ como uma combina\'c\~ao linear de $\u_1$ e
        $\u_2$. No geral, isso significa que qualquer vetor no span de $\{\a_1, \a_2\}$
        pode ser expresso na base ortogonal $\{\u_1, \u_2\}$.
        
    \end{solution}
    
    \item Mostre por indu\'c\~ao que para quaisquer $k \le n$ vetores linearmente independentes
    $\mathbf{a}_1,\ldots,\mathbf{a}_k$, os vetores $\mathbf{u}_i$, $i=1, \ldots
    k$, s\~ao ortogonais, onde
    \begin{align}
        \u_i &= \a_i - \sum_{j=1}^{i-1} \frac{\u_j^\top \a_i}{\u_j^\top\u_j} \u_j. \label{eq:Gram-Schmidt-basis-def}
    \end{align}
    O c\'alculo dos vetores $\u_i$ \'e chamado de ortogonaliza\'c\~ao de Gram–Schmidt.
    \begin{solution}
        Mostramos acima que a afirma\'c\~ao \'e v\'alida para dois vetores. Este \'e o caso
        base para a prova por indu\'c\~ao. Assuma agora que a afirma\'c\~ao seja v\'alida para $k$
        vetores. O passo indutivo na prova por indu\'c\~ao consiste ent\~ao em
        mostrar que a afirma\'c\~ao tamb\'em \'e v\'alida para $k+1$ vetores.
        
        Assuma que $\u_1, \u_2, \ldots, \u_k$ s\~ao vetores ortogonais. A hip\'otese de independ\^encia linear garante que nenhum dos
        $\u_i$ seja um vetor nulo. Temos ent\~ao para $\u_{k+1}$
        \begin{align}
            \u_{k+1} &= \a_{k+1} - \frac{\u_1^\top\a_{k+1}}{\u_1^\top\u_1}\u_1 - \frac{\u_2^\top\a_{k+1}}{\u_2^\top\u_2}\u_2
            - \ldots - \frac{\u_k^\top\a_{k+1}}{\u_k^\top\u_k}\u_k,
            \label{eq:u_k+1}
        \end{align}
        e para todo $i = 1, 2, \ldots, k$
        \begin{align}
            \u_i^\top\u_{k+1} &= \u_i^\top\a_{k+1} - \frac{\u_1^\top\a_{k+1}}{\u_1^\top\u_1}\u_i^\top\u_1 - 
            \ldots - \frac{\u_k^\top\a_{k+1}}{\u_k^\top\u_k}\u_i^\top\u_k.
        \end{align}
        Por hip\'otese $\u_i^\top\u_j = 0$ se $i \neq j$, de modo que
        \begin{align}
            \u_i^\top\u_{k+1} &= \u_i^\top\a_{k+1} - 0 - \ldots - \frac{\u_i^\top\a_{k+1}}{\u_i^\top\u_i}\u_i^\top\u_i - 0 \ldots - 0 \\
            &= \u_i^\top\a_{k+1} - \u_i^\top\a_{k+1} \\
            &= 0,
        \end{align}
        o que significa que $\u_{k+1}$ \'e ortogonal a $\u_1, \ldots, \u_k$.
        
    \end{solution}
    
    \item Mostre por indu\'c\~ao que qualquer combina\'c\~ao linear de $\a_1, \a_2, \ldots, \a_k$ (linearmente independentes)
    pode ser escrita em termos de $\u_1, \u_2, \ldots, \u_k$.
    
    \begin{solution} O caso base de dois vetores foi provado acima. Usando
        indu\'c\~ao, assumimos que a afirma\'c\~ao seja v\'alida para $k$ vetores e provaremos que
        ela tamb\'em ser\'a v\'alida para $k + 1$ vetores: Seja $\v$ uma combina\'c\~ao linear de
        $\a_1, \a_2, \ldots,
        \a_{k+1}$, i.e. $\v = \alpha_1\a_1 + \alpha_2\a_2 + \ldots + \alpha_k\a_k +
        \alpha_{k+1}\a_{k+1}$ para alguns n\'umeros reais $\alpha_1, \alpha_2, \ldots,
        \alpha_{k+1}$. Usando a hip\'otese de indu\'c\~ao, $\v$ pode ser escrito como
        \begin{equation}
            \v = \beta_1\u_1 + \beta_2\u_2 + \ldots + \beta_k\u_k + \alpha_{k+1}\a_{k+1},
        \end{equation}
        para alguns n\'umeros reais $\beta_1, \beta_2, \ldots, \beta_k$. Al\'em disso, usando a equa\'c\~ao \eqref{eq:u_k+1}, $\v$ pode ser escrito como
        \begin{align}
            \v &= \beta_1\u_1  + \ldots + \beta_k\u_k + \alpha_{k+1}\u_{k+1} + \alpha_{k+1}\frac{\u_1^\top\a_{k+1}}{\u_1^\top\u_1}\u_1 \\ 
            &\phantom{=}+ \ldots + \alpha_{k+1}\frac{\u_k^\top\a_{k+1}}{\u_k^\top\u_k}\u_k.
        \end{align}
        Com $\gamma_i = \beta_i + \alpha_{k+1}(\u_i^\top\a_{k+1})/(\u_i^\top\u_i)$, $\v$ pode, portanto, ser escrito como
        \begin{align}
            \v &= \gamma_1\u_1 + \gamma_2\u_2 + \ldots + \gamma_k\u_k + \alpha_{k+1}\u_{k+1},
        \end{align}
        o que completa a prova. No geral, isso significa que os $\u_1, \u_2, \ldots,
        \u_k$ formam uma base ortogonal para $\text{span}(\a_1, \ldots, \a_k)$, i.e.
        o conjunto de todos os vetores que podem ser obtidos pela combina\'c\~ao linear dos
        $\a_i$.
        
    \end{solution}
    
    \item Considere o caso onde $\a_1, \a_2, \ldots, \a_k$ s\~ao linearmente
    independentes e $\a_{k+1}$ \'e uma combina\'c\~ao linear de $\a_1, \a_2, \ldots,
    \a_k$. Mostre que $\u_{k+1}$, calculado de acordo com
    \eqref{eq:Gram-Schmidt-basis-def}, \'e zero.
    
    \begin{solution}
        Come\'cando com \eqref{eq:Gram-Schmidt-basis-def}, temos
        \begin{equation}
            \u_{k+1} = \a_{k+1} - \sum_{j=1}^{k} \frac{\u_j^\top \a_{k+1}}{\u_j^\top\u_j} \u_j.
        \end{equation}
        Por hip\'otese, $\a_{k+1}$ \'e uma combina\'c\~ao linear de $\a_1, \a_2, \ldots,
        \a_k$. Pela quest\~ao anterior, ele pode, portanto, tamb\'em ser escrito como uma combina\'c\~ao
        linear dos $\u_1, \ldots, \u_k$. Isso significa que existem alguns
        $\beta_i$ tais que
        \begin{equation}
            \a_{k+1} = \sum_{i=1}^k \beta_i \u_i
        \end{equation}
        se sustenta. Inserindo esta expans\~ao na equa\'c\~ao acima, obtemos
        \begin{align}
            \u_{k+1}  &=  \sum_{i=1}^k \beta_i \u_i - \sum_{j=1}^k \sum_{i=1}^k \beta_i \frac{\u_j^\top\u_i}{\u_j^\top\u_j} \u_j\\
            &= \sum_{i=1}^k \beta_i \u_i - \sum_{i=1}^k \beta_i  \frac{\u_i^\top\u_i}{\u_i^\top\u_i} \u_i
        \end{align}
        porque $\u_j^\top \u_i = 0$ se $i \neq j$. Obtemos assim o resultado desejado:
        \begin{align}
            \u_{k+1} & =  \sum_{i=1}^k \beta_i \u_i - \sum_{i=1}^k \beta_i \u_i\\
            &= 0
        \end{align}
        Esta propriedade do processo de Gram-Schmidt em
        \eqref{eq:Gram-Schmidt-basis-def} pode ser usada para verificar se uma lista de
        vetores $\a_1, \a_2, \ldots, \a_d$ \'e linearmente independente ou n\~ao. Se, por
        exemplo, $\u_{k+1}$ for zero, $\a_{k+1}$ \'e uma combina\'c\~ao linear dos $\a_1,
        \ldots, \a_k$. Al\'em disso, o resultado pode ser usado para extrair uma sublista de vetores linearmente
        independentes: remover\'iamos $\a_{k+1}$ da lista e reiniciar\'iamos
        o procedimento em \eqref{eq:Gram-Schmidt-basis-def} com $\a_{k+2}$ tomando
        o lugar de $\a_{k+1}$. Continuar desta forma constr\'oi uma lista de
        $\a_j$ linearmente independentes e $\u_j$ ortogonais, $j=1, \ldots, r$, onde
        $r$ \'e o n\'umero de vetores linearmente independentes entre os $\a_1, \a_2,
        \ldots, \a_d$.
        
    \end{solution}
\end{exenumerate}

% --------------------------------------------------
\ex{Transforma\'c\~oes lineares}
\label{ex:linear-transforms}

\begin{exenumerate}
    \item Assuma que dois vetores $ \a_1$ e $\a_2$ est\~ao em $\mathbb{R}^2$. Juntos,
    eles definem um paralelogramo. Use a \exref{ex:gram-schmidt} para mostrar que a
    \'area ao quadrado $S^2$ do paralelogramo \'e dada por
    \begin{equation}
        S^2 = (\a_2^\top\a_2)(\a_1^\top\a_1)-(\a_2^\top\a_1)^2
    \end{equation}
    
    \begin{solution}
        Sejam $\a_1$ e $\a_2$ os vetores que definem o paralelogramo. Da geometria,
        sabemos que a \'area de um paralelogramo \'e base vezes altura, o que \'e
        equivalente ao comprimento do vetor da base vezes o comprimento do vetor da altura.
        Denotamos isso por $S^2 = ||\a_1||^2||\u_2||^2$, onde $\a_1$ \'e o vetor da base
        e $\u_2$ \'e o vetor da altura que \'e ortogonal ao vetor da base. Usando o processo de
        Gram--Schmidt para os vetores $\a_1$ e $\a_2$ nessa
        ordem, obtemos o vetor $\u_2$ como a segunda sa\'ida.
        % 
        \begin{figure}[h!]
            \centering
            \scalebox{0.75}{
                \input{figs/linear-algebra/gram.pstex_t} }
        \end{figure}
        
        Portanto $||\u_2||^2$ \'e igual a
        \begin{align}
            ||\u_2||^2 &= \u_2^\top\u_2 \\
            &= \left(\a_2 - \frac{\a_1^\top\a_2}{\a_1^\top\a_1}\a_1\right)^\top\left(\a_2 - \frac{\a_1^\top\a_2}{\a_1^\top\a_1}\a_1\right) \\
            &= \a_2^\top\a_2 - \frac{(\a_1^\top\a_2)^2}{\a_1^\top\a_1} - \frac{(\a_1^\top\a_2)^2}{\a_1^\top\a_1} + \left(\frac{\a_1^\top\a_2}{\a_1^\top\a_1}\right)^2\a_1^\top\a_1 \\
            &= \a_2^\top\a_2 - \frac{(\a_1^\top\a_2)^2}{\a_1^\top\a_1}.
        \end{align}
        Assim, $S^2$ \'e:
        \begin{align}
            S^2 &= ||\a_1||^2||\u_2||^2 \\
            &= (\a_1^\top\a_1)(\u_2^\top\u_2) \\
            &= (\a_1^\top\a_1)\left(\a_2^\top\a_2 - \frac{(\a_1^\top\a_2)^2}{\a_1^\top\a_1}\right) \\
            &= (\a_2^\top\a_2)(\a_1^\top\a_1) - (\a_1^\top\a_2)^2.
        \end{align}
        
    \end{solution}
    
    \item Forme a matriz $\A=(\a_1 \; \a_2)$ onde $\a_1$ e $\a_2$ s\~ao
    o primeiro e o segundo vetor coluna, respectivamente. Mostre que
    \begin{equation}
        S^2 = (\det \A)^2.
    \end{equation}
    
    
    \begin{solution}
        Formamos a matriz $\A$,
        \begin{equation}
            \A = \begin{pmatrix} \a_1 & \a_2
            \end{pmatrix}  =
            \begin{pmatrix} a_{11} & a_{12} \\ a_{21} & a_{22}
            \end{pmatrix}.
        \end{equation}
        O determinante de $\A$ \'e $\det \A = a_{11}a_{22} - a_{12}a_{21}$. Ao expandir $(\a_2^\top\a_2)$, $(\a_1^\top\a_1)$ e 
        $(\a_1^\top\a_2)^2$, obtemos
        \begin{align}
            \a_2^\top\a_2   &= a_{12}^2 + a_{22}^2 \\ 
            \a_1^\top\a_1   &= a_{11}^2 + a_{21}^2 \\
            (\a_1^\top\a_2)^2   &= (a_{11}a_{12} + a_{21}a_{22})^2 = a_{11}^2a_{12}^2 + a_{21}^2a_{22}^2 + 2a_{11}a_{12}a_{21}a_{22}.
        \end{align}
        Portanto a \'area \'e igual a
        \begin{align}
            S^2 &= (a_{12}^2 + a_{22}^2)(a_{11}^2 + a_{21}^2) - (\a_1^\top\a_2)^2 \\
            &= a_{12}^2a_{11}^2 + a_{12}^2a_{21}^2 + a_{22}^2a_{11}^2 + a_{22}^2a_{21}^2 \nonumber\\
            &\phantom{=}- (a_{12}^2a_{11}^2 + a_{21}^2a_{22}^2 + 2a_{11}a_{12}a_{21}a_{22}) \\
            &= a_{12}^2a_{21}^2 + a_{22}^2a_{11}^2 - 2a_{11}a_{12}a_{21}a_{22} \\
            &= (a_{11}a_{22} - a_{12}a_{21})^2,
        \end{align}
        que \'e igual a $(\det \A)^2$.
    \end{solution}
    \item Considere a transforma\'c\~ao linear $\y=\A\x$ onde $\A$
    \'e uma matriz $2 \times 2$. Denomine a imagem do ret\^angulo $U_x=[x_1 \;
    x_1+\triangle_1]\times [x_2 \; x_2+\triangle_2]$ sob a transforma\'c\~ao
    $\A$ por $U_y$. O que \'e $U_y$? Qual \'e a \'area de $U_y$?
    
    \begin{solution} $U_y$ \'e um paralelogramo que \'e definido pelos vetores
        coluna $\a_1$ e $\a_2$ de $\A$, quando $\A = (\a_1 \; \a_2)$.
        
        Um ret\^angulo com a mesma \'area que $U_x$ \'e definido pelos vetores $(\Delta_1,
        0)$ e $(0, \Delta_2)$. Sob a transforma\'c\~ao linear $\A$, esses vetores de base
        tornam-se $\Delta_1\a_1$ e $\Delta_2\a_2$. Portanto, um paralelogramo
        com a mesma \'area que $U_y$ \'e definido por $\Delta_1\a_1$ e
        $\Delta_2\a_2$, como mostrado na figura seguinte.
        \begin{figure}[htp] \centering \scalebox{0.4}{
                \begin{picture}(0,0)%
                    \includegraphics{linear1.pdf}%
                \end{picture}%
                \setlength{\unitlength}{4144sp}%
                % 
                \begingroup\makeatletter\ifx\SetFigFont\undefined%
                \gdef\SetFigFont#1#2#3#4#5{%
                    \reset@font\fontsize{#1}{#2pt}%
                    \fontfamily{#3}\fontseries{#4}\fontshape{#5}%
                    \selectfont}%
                \fi\endgroup%
                \begin{picture}(6417,4299)(-1454,152)
                    \put(2026,3809){\makebox(0,0)[lb]{\smash{{\SetFigFont{25}{30.0}{\familydefault}{\mddefault}{\updefault}{\color[rgb]{0,0,0}$U_x$}%
                    }}}}
                    \put(-629,1649){\makebox(0,0)[lb]{\smash{{\SetFigFont{25}{30.0}{\familydefault}{\mddefault}{\updefault}{\color[rgb]{0,0,0}$x_2$}%
                    }}}}
                    \put(3016,434){\makebox(0,0)[lb]{\smash{{\SetFigFont{25}{30.0}{\familydefault}{\mddefault}{\updefault}{\color[rgb]{0,0,0}$x_1 + \Delta_1$}%
                    }}}}
                    \put(-1439,3449){\makebox(0,0)[lb]{\smash{{\SetFigFont{25}{30.0}{\familydefault}{\mddefault}{\updefault}{\color[rgb]{0,0,0}$x_2 + \Delta_2$}%
                    }}}}
                    \put(721,434){\makebox(0,0)[lb]{\smash{{\SetFigFont{25}{30.0}{\familydefault}{\mddefault}{\updefault}{\color[rgb]{0,0,0}$x_1$}%
                    }}}}
                \end{picture}%
            }
            % 
            \scalebox{0.4}{ 
                \begin{picture}(0,0)%
                    \includegraphics{linear2.pdf}%
                \end{picture}%
                \setlength{\unitlength}{4144sp}%
                % 
                \begingroup\makeatletter\ifx\SetFigFont\undefined%
                \gdef\SetFigFont#1#2#3#4#5{%
                    \reset@font\fontsize{#1}{#2pt}%
                    \fontfamily{#3}\fontseries{#4}\fontshape{#5}%
                    \selectfont}%
                \fi\endgroup%
                \begin{picture}(5649,4299)(-686,152)
                    \put(1711,1199){\makebox(0,0)[lb]{\smash{{\SetFigFont{17}{20.4}{\familydefault}{\mddefault}{\updefault}{\color[rgb]{0,0,0}$x_1$}%
                    }}}}
                    \put(3241,1694){\makebox(0,0)[lb]{\smash{{\SetFigFont{17}{20.4}{\familydefault}{\mddefault}{\updefault}{\color[rgb]{0,0,0}$x_1 + \Delta_1$}%
                    }}}}
                    \put( 91,3044){\makebox(0,0)[lb]{\smash{{\SetFigFont{17}{20.4}{\familydefault}{\mddefault}{\updefault}{\color[rgb]{0,0,0}$x_2 + \Delta_2$}%
                    }}}}
                    \put(181,1874){\makebox(0,0)[lb]{\smash{{\SetFigFont{17}{20.4}{\familydefault}{\mddefault}{\updefault}{\color[rgb]{0,0,0}$x_2$}%
                    }}}}
                    \put(991,929){\makebox(0,0)[lb]{\smash{{\SetFigFont{20}{24.0}{\familydefault}{\mddefault}{\updefault}{\color[rgb]{0,0,0}$\mathbf{a}_1$}%
                    }}}} 
                    \put(406,1289){\makebox(0,0)[lb]{\smash{{\SetFigFont{20}{24.0}{\familydefault}{\mddefault}{\updefault}{\color[rgb]{0,0,0}$\mathbf{a}_2$}%
                    }}}}
                    \put(3691,4214){\makebox(0,0)[lb]{\smash{{\SetFigFont{25}{30.0}{\familydefault}{\mddefault}{\updefault}{\color[rgb]{0,0,0}$U_y$}%
                    }}}}
                \end{picture}%
            }
        \end{figure}
        
        Pela quest\~ao anterior, a $A_{U_y}$ de $U_y$ \'e igual ao valor absoluto do determinante da matriz $(\Delta_1\a_1 \; \Delta_2\a_2)$:
        \begin{align}
            A_{U_y} &= |\det \begin{pmatrix} \Delta_1a_{11} & \Delta_2a_{12} \\ \Delta_1a_{21} & \Delta_2a_{22} \end{pmatrix} |\\
            &= |\Delta_1\Delta_2a_{11}a_{22} - \Delta_1\Delta_2a_{12}a_{21}| \\ 
            &= |\Delta_1\Delta_2(a_{11}a_{22} - a_{12}a_{21})| \\
            &= \Delta_1\Delta_2 |\det \A|
        \end{align}
        Portanto, a \'area de $U_y$ \'e a \'area de $U_x$ vezes $|\det \A|$.
    \end{solution}
    
    \item Forne\'ca uma explica\'c\~ao intuitiva do porqu\^e temos igualdade na f\'ormula de mudan\'ca de vari\'aveis 
    \begin{equation} \int_{U_y} f(\y)d\y = \int_{U_x} f(\A \x) |\det \A| d\x.
    \end{equation} onde $\A$ \'e tal que $U_x$ \'e um (hiper-)ret\^angulo alinhado aos eixos como na quest\~ao anterior.
    
    \begin{solution} Podemos pensar que, de modo informal, as duas integrais
        s\~ao limites das duas somas seguintes
        \begin{equation} \sum_{\y_i \in U_y}
            f(\y_i) \textrm{vol}(\Delta_{\y_i}) \quad \quad \quad \sum_{\x_i \in U_x}
            f(\A \x_i) |\det \A| \textrm{vol}(\Delta_{\x_i})
        \end{equation}
        onde $\x_i = \A^{-1} \y_i$, o que significa que $\x$ e $\y$ est\~ao relacionados
        por $\y = \A \x$. O conjunto de valores da fun\'c\~ao $f(\y_i)$ e $f(\A \x_i)$
        que entram nas duas somas s\~ao exatamente os mesmos. O volume
        $\textrm{vol}(\Delta_{\x_i})$ de um pequeno hipercubo alinhado aos eixos (em $d$
        dimens\~oes) \'e igual a $\prod_{i=1}^d \Delta_i$. A imagem deste pequeno
        hipercubo alinhado aos eixos sob $\A$ \'e um paralelogramo $\Delta_{\y_i}$
        com volume $\textrm{vol}(\Delta_{\y_i})= |\det \A|
        \textrm{vol}(\Delta_{\x_i})$. Assim
        \begin{equation}
            \sum_{\y_i \in U_y} f(\y_i) \textrm{vol}(\Delta_{\y_i}) = \sum_{\x_i \in U_x} f(\A \x_i) |\det \A| \textrm{vol}(\Delta_{\x_i}).
        \end{equation}
        Devemos ter o termo $|\det \A|$ para compensar o fato de que o volume de $U_x$ e $U_y$ n\~ao s\~ao os mesmos. Por exemplo, seja $\A$ uma matriz diagonal $\diag(10, 100)$ de modo que $U_x$ seja muito menor que $U_y$. O determinante $\det \A = 1000$ ent\~ao compensa o fato de que os valores de $\x_i$ est\~ao mais condensados do que os de $\y_i$.
    \end{solution}
    
\end{exenumerate}


% --------------------------------------------------

\ex{Decomposi\'c\~ao em autovalores}
\label{ex:eigenvalue-decomposition}
Para uma matriz quadrada $\A$ de tamanho $n \times n$, um vetor $\u_i
\neq 0$ que satisfaz
\begin{equation}
    \A\u_i = \lambda_i \u_i
    \label{eq:eigenvector}
\end{equation}
\'e chamado de autovetor de $\A$, e $\lambda_i$ \'e o autovalor correspondente.
Para uma matriz de tamanho $n \times n$, existem $n$
autovalores $\lambda_i$ (que n\~ao s\~ao necessariamente distintos).
\begin{exenumerate}
    \item Mostre que se $\u_1$ e $\u_2$ s\~ao autovetores com $\lambda_1=\lambda_2$, ent\~ao
    $\u = \alpha \u_1 +\beta \u_2$ tamb\'em \'e um autovetor com o mesmo autovalor.
    \begin{solution}
        Calculamos
        \begin{align}
            \A\u &= \alpha \A\u_1 + \beta \A\u_2 \\
            &= \alpha \lambda \u_1 + \beta \lambda \u_2 \\
            &= \lambda(\alpha\u_1 + \beta\u_2) \\
            &= \lambda \u,
        \end{align}
        ent\~ao $\u$ \'e um autovetor de $\A$ com o mesmo autovalor que $\u_1$ e $\u_2$.
    \end{solution}
    
    \item Assuma que nenhum dos autovalores de $\A$ seja zero. Denote por $\U$
    a matriz onde os vetores coluna s\~ao autovetores linearmente independentes
    $\u_i$ de $\A$. Verifique que \eqref{eq:eigenvector} pode ser escrito em forma de matriz como $\A\U=\U\Lambdab$, onde $\Lambdab$ \'e uma
    matriz diagonal com os autovalores $\lambda_i$ como elementos da diagonal.
    
    \begin{solution}
        Pelas propriedades b\'asicas da multiplica\'c\~ao de matrizes, temos
        \begin{align}
            \A\U &= (\A\u_1 \; \A\u_2 \;\ldots \; \A\u_n)
        \end{align}
        Com $\A\u_i = \lambda_i \u_i$ para todo $i = 1, 2, \ldots, n$, obtemos assim
        \begin{align}
            \A\U & = (\lambda_1\u_1 \; \lambda_2\u_2 \;\ldots \; \lambda_n\u_n) \\
            &= \U\Lambdab.
        \end{align}
    \end{solution}
    
    \item Mostre que podemos escrever, com $\V^\top=\U^{-1}$,
    \begin{align}
        \A &= \U \Lambdab \V^\top,&     \A &= \sum_{i=1}^n \lambda_i \u_i\v_i^\top,\\
        \A^{-1} &= \U \Lambdab^{-1} \V^\top, & \A^{-1} &= \sum_{i=1}^n \frac{1}{\lambda_i} \u_i\v_i^\top,
    \end{align}
    onde $\v_i$ \'e a $i$-\'esima coluna de $\V$.
    \begin{solution}
        \begin{enumerate}
            \item[(i)] Como as colunas de $\U$ s\~ao linearmente independentes,
            $\U$ \'e invers\'ivel. Como $\A\U = \U\Lambdab$, multiplicar pela
            direita pela inversa de $\U$ resulta em $\A = \U\Lambdab \U^{-1} =
            \U\Lambda \V^\top$.
            \item[(ii)] Denomine por $\u^{[i]}$ a $i$-\'esima linha de $\U$,
            $\v^{(j)}$ a $j$-\'esima coluna de $\V^\top$ e $\v^{[j]}$
            a $j$-\'esima linha de $\V$, e denote $\B = \sum_{i=1}^n
            \lambda_i\u_i\v_i^\top$. Seja $\e^{[i]}$ um vetor linha
            com 1 na $i$-\'esima posi\'c\~ao e 0 nos outros lugares, e $\e^{(j)}$ um
            vetor coluna com 1 na $j$-\'esima posi\'c\~ao e 0
            nos outros lugares. Note que porque $\A = \U\Lambdab \V^\top$,
            o elemento na $i$-\'esima linha e $j$-\'esima coluna \'e
            \begin{align}
                A_{ij}  &= \u^{[i]}\Lambdab \v^{(j)} \\
                &= \u^{[i]}\Lambdab {\v^{[j]}}^\top \\
                &= \u^{[i]} \left( \begin{matrix} \lambda_1V_{j1} \\ \vdots \\ \lambda_nV_{jn} \end{matrix} \right) \\
                &= \sum_{k=1}^n \lambda_k V_{jk}U_{ik}.
            \end{align}
            Por outro lado, para a matriz $\B$, o elemento na $i$-\'esima linha e $j$-\'esima coluna \'e
            \begin{align}
                B_{ij}  &= \sum_{k=1}^n \lambda_k\e^{[i]}\u_k\v_k^\top \e^{(j)} \\
                &= \sum_{k=1}^n \lambda_kU_{ik}V_{jk},
            \end{align}
            que \'e o mesmo que $A_{ij}$. Portanto $\A = \B$.
            
            \item[(iii)] Como $\Lambdab$ \'e uma matriz diagonal sem zeros como elementos da diagonal,
            ela \'e invers\'ivel. Temos, portanto
            \begin{align}
                \A^{-1} &= (\U\Lambdab{\U^{-1}})^{-1} \\
                & = (\Lambdab \U^{-1})^{-1}\U^{-1} \\
                &= \U \Lambda^{-1} \U^{-1} \\
                &= \U\Lambda^{-1}\V^\top.
            \end{align}
            
            \item[(iv)] Isso decorre de $\A = \U \Lambdab \V^\top = \sum_i \u_i \lambda_i \v_i^\top$, quando $\lambda_i$ \'e substitu\'ido por $1/\lambda_i$.
        \end{enumerate}
    \end{solution}  
\end{exenumerate}

% --------------------------------------------------

\ex{Tra\'co, determinantes e autovalores}
\label{ex:trace-determinants-eigenvalues}
\begin{exenumerate}
    \item Use a \exref{ex:eigenvalue-decomposition} para mostrar que $\tr(\A)=\sum_i A_{ii} = \sum_i\lambda_i$. (Voc\^e pode usar $\tr(\A\B)=\tr(\B\A)$.)
    \begin{solution}
        Como $\tr(\A\B) = \tr(\B\A)$ e $\A = \U\Lambdab \U^{-1}$
        \begin{align}
            \tr(\A)  &= \tr(\U\Lambdab \U^{-1}) \\
            &= \tr(\Lambdab \U^{-1}\U) \\
            &= \tr(\Lambda)\\
            &= \sum_i \lambda_i.
        \end{align}
    \end{solution}
    
    \item Use a \exref{ex:eigenvalue-decomposition} para mostrar que $\det \A = \prod_i \lambda_i$. (Use $\det \A^{-1}=1/(\det \A)$ e $\det(\A\B)=\det(\A)\det(\B)$ para quaisquer $\A$ e $\B$.)
    \begin{solution}
        Utilizamos a decomposi\'c\~ao em autovalores de $\A$ para obter
        \begin{align}
            \det(\A) &= \det(\U\Lambdab \U^{-1})\\
            &= \det(\U)\det(\Lambdab)\det(\U^{-1}) \\
            &= \frac{\det(\U)\det(\Lambdab)}{\det(\U)}\\
            &= \det(\Lambda)\\
            &= \prod_i \lambda_i,
        \end{align}
        onde, na \'ultima linha, usamos que o determinante de uma matriz diagonal \'e o produto de seus elementos da diagonal.
        
    \end{solution}
\end{exenumerate}


% --------------------------------------------------

\ex{Decomposi\'c\~ao em autovalores para matrizes sim\'etricas}
\label{ex:eigenvalue-decomposition-symmetric-matrices}
\begin{exenumerate}
    \item Assuma que uma matriz $\A$ seja sim\'etrica, i.e. $\A^\top=\A$. Sejam $\u_1$ e
    $\u_2$ dois autovetores de $\A$ com autovalores correspondentes $\lambda_1$
    e $\lambda_2$, com $\lambda_1 \neq \lambda_2$. Mostre que os dois vetores
    s\~ao ortogonais entre si.
    
    \begin{solution}
        Como $\A\u_2 = \lambda_2\u_2$, temos
        \begin{equation}
            \u_1^\top \A\u_2 = \lambda_2\u_1^\top\u_2.
        \end{equation}
        Tomando a transposta de $\u_1^\top \A\u_2$ obtemos
        \begin{align}
            (\u_1^\top \A\u_2)^\top &= (\A\u_2)^\top(\u_1^\top)^\top = \u_2^\top \A^\top \u_1 = \u_2^\top \A\u_1 \\
            & = \lambda_1\u_2^\top\u_1
        \end{align}
        porque $\A$ \'e sim\'etrica e $\A \u_1 = \lambda_1 \u_1$. Por outro lado, a mesma opera\'c\~ao resulta em
        \begin{align}
            (\u_1^\top \A\u_2)^\top &= (\lambda_2\u_1^\top\u_2)^\top = \lambda_2\u_2^\top\u_1
        \end{align}
        Portanto $\lambda_1\u_2^\top\u_1 = \lambda_2\u_2^\top\u_1$, o que \'e equivalente a $\u_2^\top\u_1(\lambda_1 - \lambda_2) = 0$. Como
        $\lambda_1 \neq \lambda_2$, a \'unica possibilidade \'e que
        $\u_2^\top\u_1 = 0$. Assim, $\u_1$ e $\u_2$ s\~ao ortogonais entre si.
        
        O resultado implica que os autovetores de uma matriz sim\'etrica $\A$ com
        autovalores distintos $\lambda_i$ formam uma base ortogonal. O resultado
        se estende ao caso onde alguns dos autovalores s\~ao iguais (n\~ao provado).
        
        
    \end{solution}
    
    \item Uma matriz sim\'etrica $\A$ \'e dita ser positiva definida se $\v^\top\A\v>0$
    para todos os vetores $\v$ n\~ao nulos. Mostre que a positividade definida implica que
    $\lambda_i>0$, $i=1,\ldots, M$. Mostre que, vice-versa, $\lambda_i>0$,
    $i=1 \ldots M$ implica que a matriz $\A$ seja positiva definida. Conclua
    que uma matriz positiva definida \'e invers\'ivel.
    
    \begin{solution}
        
        Assuma que $\v^\top \A\v > 0$ para todo $\v \neq 0$. Como autovetores n\~ao s\~ao
        vetores nulos, a suposi\'c\~ao tamb\'em vale para o autovetor $\u_k$ com
        autovalor correspondente $\lambda_k$. Agora
        \begin{align}
            \u_k^\top \A\u_k &= \u_k^\top\lambda_k\u_k = \lambda_k(\u_k^\top\u_k) = \lambda_k||\u_k|| > 0
        \end{align}
        e porque $||\u_k|| > 0$, obtemos $\lambda_k > 0$.
        
        Assuma agora que todos os autovalores de $\A$, $\lambda_1, \lambda_2, \ldots, \lambda_n$, s\~ao positivos e n\~ao nulos. Mostramos acima que existe uma base ortogonal consistindo de autovetores $\u_1, \u_2, \ldots, \u_n$ e, portanto, todo vetor $\v$ pode ser escrito como uma
        combina\'c\~ao linear desses vetores (apenas mostramos para o caso de autovalores distintos, mas vale de forma geral). Assim, para um vetor n\~ao nulo $\v$ e para alguns n\'umeros reais $\alpha_1, \alpha_2, \ldots, \alpha_n$, temos
        \begin{align}
            \v^\top\A\v &= (\alpha_1\u_1 + \ldots + \alpha_n\u_n)^\top \A(\alpha_1\u_1 + \ldots + \alpha_n\u_n) \\
            &= (\alpha_1\u_1 + \ldots + \alpha_n\u_n)^\top(\alpha_1\A\u_1 + \ldots + \alpha_n\A\u_n) \\
            &= (\alpha_1\u_1 + \ldots + \alpha_n\u_n)^\top(\alpha_1\lambda_1\u_1 + \ldots + \alpha_n\lambda_n\u_n) \\
            &= \sum_{i,j} \alpha_i\u_i^\top\alpha_j\lambda_j\u_j \\
            &= \sum_i \alpha_i\alpha_i\lambda_i\u_i^\top\u_i \\
            &= \sum_i (\alpha_i)^2||\u_i||^2\lambda_i,
        \end{align}
        onde usamos que $\u_i^\top\u_j = 0$ se $i \neq j$, devido \`a ortogonalidade da base. Como $(\alpha_i)^2 > 0$, $||\u_i||^2 > 0$ e $\lambda_i > 0$ para todo $i$, descobrimos que $\v^\top\A\v > 0.$
        
        Como todo autovalor de $\A$ \'e n\~ao nulo, podemos usar a \exref{ex:eigenvalue-decomposition} para concluir que a inversa de $\A$ existe e \'e igual a $\sum_i 1/\lambda_i \u_i \u_i^\top$.
        
    \end{solution}
    
\end{exenumerate}


% --------------------------------------


\ex{M\'etodo das pot\^encias}
\label{ex:power-method}
Analisamos aqui um algoritmo chamado de ``m\'etodo das pot\^encias''. O m\'etodo das pot\^encias toma
como entrada uma matriz sim\'etrica positiva definida $\Sigmab$ e calcula o
autovetor que possui o maior autovalor (o ``primeiro autovetor''). Por
exemplo, no caso da an\'alise de componentes principais, $\Sigmab$ \'e a matriz de covari\^ancia dos dados observados e o primeiro autovetor \'e a dire\'c\~ao do primeiro componente principal.

O m\'etodo das pot\^encias consiste em iterar as equa\'c\~oes de atualiza\'c\~ao
\begin{align}
    \v_{k+1} &= \Sigmab \w_{k}, & \w_{k+1} &= \frac{\v_{k+1}}{||\v_{k+1}||_2},
\end{align}
onde $||\v_{k+1}||_2$ denota a norma Euclidiana.

\begin{exenumerate}
    \item Seja $\U$ a matriz com os autovetores (ortonormais) $\u_i$ de $\Sigmab$ como colunas. Qual \'e a decomposi\'c\~ao em autovalores da matriz de covari\^ancia $\Sigmab$?
    
    \begin{solution}
        Como as colunas de $\U$ s\~ao vetores (autovetores) ortonormais, $\U$ \'e ortogonal, i.e.
        $\U^{-1} = \U^\top$. Com a \exref{ex:eigenvalue-decomposition} e a \exref{ex:eigenvalue-decomposition-symmetric-matrices}, obtemos
        \begin{align}
            \Sigmab &= \U\Lambdab \U^\top,
        \end{align}
        onde $\Lambdab$ \'e a matriz diagonal com autovalores $\lambda_i$ de $\Sigmab$ como elementos da diagonal. Deixe os autovalores ordenados como $\lambda_1 > \lambda_2 > \ldots > \lambda_n > 0$ (e, como hip\'otese adicional, todos distintos).
    \end{solution}
    
    
    \item Sejam $\tilde{\v}_{k}= \U^\top \v_k$ e $\tilde{\w}_{k}= \U^\top \w_k$. Escreva as equa\'c\~ao de atualiza\'c\~ao do m\'etodo das pot\^encias em termos de $\tilde{\v}_{k}$ e $\tilde{\w}_{k}$. Isso significa que estamos fazendo uma mudan\'ca de base para representar os vetores $\w_k$ e $\v_k$ na base dada pelos autovetores de $\Sigmab$.
    
    \begin{solution}
        Com
        \begin{align}
            \v_{k+1} &= \Sigmab \w_k \\
            & = \U\Lambdab \U^\top\w_k
        \end{align}
        obtemos
        \begin{align}
            \U^\top\v_{k+1} = \Lambdab \U^\top\w_k.
        \end{align}
        Assim $\tilde{\v}_{k+1} = \Lambdab\tilde{\w}_k$. A norma de $\tilde{\v}_{k+1}$ \'e a mesma que a norma de $\v_{k+1}$:
        \begin{align}
            ||\tilde{\v}_{k+1}||_2 &= ||\U^\top\v_{k+1}||_2 \\
            &= \sqrt{(\U^\top\v_{k+1})^\top(\U^\top\v_{k+1})} \\
            &= \sqrt{\v_{k+1}^\top \U\U^\top\v_{k+1}} \\
            &= \sqrt{\v_{k+1}^\top\v_{k+1}} \\
            &= ||\v_{k+1}||_2.
        \end{align}
        Assim, a equa\'c\~ao de atualiza\'c\~ao, em termos de $\tilde{\v}_{k}$ e $\tilde{\w}_k$, \'e
        \begin{align}
            \tilde{\v}_{k+1} &= \Lambdab \tilde{\w}_k, & \tilde{\w}_{k+1} &= \frac{\tilde{\v}_{k+1}}{||\tilde{\v}_{k+1}||}.
        \end{align}
        
    \end{solution}
    
    
    \item Assuma que voc\^e comece a itera\'c\~ao com $\tilde{\w}_{0}$. Para qual vetor $\tilde{\w}^{\ast}$ a itera\'c\~ao converge? 
    
    \begin{solution}
        Seja $\tilde{\w}_0 = \begin{pmatrix} \alpha_1 & \alpha_2 & \ldots & \alpha_n \end{pmatrix}^\top$. Como $\Lambdab$ \'e uma matriz diagonal, obtemos
        \begin{align}
            \tilde{\v}_1 &=  
            \begin{pmatrix} 
                \lambda_1\alpha_1 \\ 
                \lambda_2\alpha_2 \\ 
                \vdots \\ 
                \lambda_n\alpha_n 
            \end{pmatrix} 
            = 
            \lambda_1\alpha_1
            \begin{pmatrix} 
                1 \\ 
                \frac{\alpha_2}{\alpha_1}\frac{\lambda_2}{\lambda_1} \\ 
                \vdots \\ 
                \frac{\alpha_n}{\alpha_1}\frac{\lambda_n}{\lambda_1}
            \end{pmatrix}
        \end{align}
        e portanto 
        \begin{align}
            \tilde{\w}_1 &= \frac{\lambda_1\alpha_1}{c_1}
            \begin{pmatrix} 1 \\
                \frac{\alpha_2}{\alpha_1}\frac{\lambda_2}{\lambda_1} \\ 
                \vdots \\ 
                \frac{\alpha_n}{\alpha_1}\frac{\lambda_n}{\lambda_1}
            \end{pmatrix},
        \end{align}
        onde $c_1$ \'e uma constante de normaliza\'c\~ao tal que $\|\tilde{\w}_1\| = 1$ (i.e. $c_1 = \|\tilde{\v}_1\|$). Logo, para $\tilde{\w}_k$ vale que
        \begin{align}
            \tilde{\w}_k &= \tilde{c}_k 
            \begin{pmatrix}
                1 \\
                \frac{\alpha_2}{\alpha_1}\left(\frac{\lambda_2}{\lambda_1}\right)^k \\
                \vdots \\
                \frac{\alpha_n}{\alpha_1}\left(\frac{\lambda_n}{\lambda_1}\right)^k 
            \end{pmatrix},
        \end{align}
        onde $\tilde{c}_k$ \'e novamente uma constante de normaliza\'c\~ao tal que $||\tilde{\w}_k|| = 1$.
        
        Como $\lambda_1$ \'e o autovalor dominante, $|\lambda_j/\lambda_1| < 1$ para $j = 2, 3, \ldots, n$, de modo que
        \begin{equation}
            \lim_{k \rightarrow \infty} \left(\frac{\lambda_j}{\lambda_1} \right)^k = 0, \quad j=2, 3, \ldots, n,
        \end{equation}
        e portanto
        \begin{equation}
            \lim_{k\to \infty}
            \begin{pmatrix}
                1 \\
                \frac{\alpha_2}{\alpha_1}\left(\frac{\lambda_2}{\lambda_1}\right)^k \\
                \vdots \\
                \frac{\alpha_n}{\alpha_1}\left(\frac{\lambda_n}{\lambda_1}\right)^k 
            \end{pmatrix}
            = \begin{pmatrix}
                1 \\
                0\\
                \vdots\\
                0
            \end{pmatrix}.
        \end{equation}
        
        Para a constante de normaliza\'c\~ao $\tilde{c}_k$, obtemos
        \begin{equation}
            \tilde{c}_k = \frac{1}{\sqrt{1 + \sum_{i=2}^n \left(\frac{\alpha_i}{\alpha_1}\right)^2 \left(\frac{\lambda_i}{\lambda_1}\right)^{2k}}},
        \end{equation}
        e portanto
        \begin{align}
            \lim_{k \rightarrow \infty} \tilde{c}_k &= \frac{1}{\sqrt{1 + \sum_{i=2}^n \left(\frac{\alpha_i}{\alpha_1}\right)^2 \underset{k \rightarrow \infty}
                    {\lim}\left(\frac{\lambda_i}{\lambda_1}\right)^{2k}}} \\
            &= \frac{1}{\sqrt{1 + \sum_{i=2}^n \left(\frac{\alpha_i}{\alpha_1}\right)^2 \cdot 0}} \\
            &= 1.
        \end{align}
        O limite do produto de duas sequ\^encias convergentes \'e o produto dos limites, de modo que 
        \begin{align}
            \lim_{k \rightarrow \infty} \tilde{\w}_k &= \lim_{k \rightarrow \infty} \tilde{c}_k 
            \lim_{k \rightarrow \infty} \begin{pmatrix}
                1 \\
                \frac{\alpha_2}{\alpha_1}\left(\frac{\lambda_2}{\lambda_1}\right)^k \\
                \vdots \\
                \frac{\alpha_n}{\alpha_1}\left(\frac{\lambda_n}{\lambda_1}\right)^k 
            \end{pmatrix}
            = 
            \begin{pmatrix} 
                1 \\ 
                0 \\ 
                \vdots \\ 
                0 
            \end{pmatrix}.
        \end{align}
        
    \end{solution}
    \item Conclua que o m\'etodo das pot\^encias encontra o primeiro autovetor.
    
    \begin{solution}
        Como $\w_k = \U\tilde{\w}_k$, obtemos
        \begin{align}
            \lim_{k \rightarrow \infty} \w_k &= \U 
            \begin{pmatrix}
                1 \\
                0 \\
                \vdots \\
                0
            \end{pmatrix}
            = \u_1,
        \end{align}
        que \'e o autovetor com o maior autovalor, i.e. o autovetor ``primeiro'' ou ``dominante''. 
    \end{solution}
    
\end{exenumerate}